\textbf{Automated Scientific Discovery.} 
Recent progress has significantly reshaped the role of AI in automating end-to-end scientific research.
AI Scientist-v1~\citep{lu2024ai} was an early milestone, showcasing how advanced language models can autonomously generate research ideas, run experiments, and draft scientific papers.
This work was followed by a series of subsequent studies in machine learning fields~\citep{zochi2025,tang2025ai} and diverse scientific disciplines~\citep{villaescusa2025denario, mitchener2025kosmos} that further advanced this line of research.
However, these approaches often suffer from an overly ambitious problem setting that aims to achieve fully automated science and tend to lack clearly defined scientific objectives for AI Scientists. Without specific goals, such systems often produce undirected discoveries that lack genuine scientific value. 
To address this issue, our Jr. AI Scientist builds on existing baselines and conducts research within a well-defined research workflow, aiming to generate higher-quality scientific papers.
As concurrent work, DeepScientist~\citep{weng2025deepscientist} also adopts a baseline-based approach. 
DeepScientist~\citep{weng2025deepscientist} focuses mainly on experimental performance, formalizing discovery as a Bayesian Optimization problem. However, its overall framework design and workflow integration are outlined at a high level, with limited discussion of implementation details. In contrast, our Jr. AI Scientist explicitly articulates each stage of the research process and further aims to contribute to the community by comprehensively reporting the failures and risks encountered throughout scientific exploration.

\textbf{AI-Assisted Scientific Research.} 
Research specialized for each element of the research process has also been actively explored~\citep{chen2025ai4research}.
For the idea generation phase, \cite{si2024can} investigates the novelty of the LLM-generated ideas, and \cite{si2025ideation} investigates the ideation–execution gap.
For the survey phase, OpenScholar~\citep{asai2024openscholar} have been developed to support literature review.
For the experimental phase, AlphaEvolve~\citep{novikov2025alphaevolve} leverages large-scale trial-and-error strategies to enhance the performance.
For the writing and review phase, CycleResearcher~\citep{weng2024cycleresearcher} provides a learning framework specialized for scientific writing, while DeepReviewer~\citep{zhu2025deepreview} focuses on the review process. Recent studies provide a comprehensive overview of automated review, outlining key challenges, proposing a practical review pipeline for real-world implementation, and constructing a large-scale dataset to support automated review research~\citep{ZHUANG2025103332, Lin2023ASPR, lin2023moprd}.
Beyond these, rather than pursuing full automation like AI Scientists, AI Co-Scientist~\citep{gottweis2025towards} emphasizes collaboration between humans and AI.
In this work, instead of focusing on individual parts of the research process, we investigate the entire end-to-end research cycle, aiming to rigorously evaluate both the performance and the limitations.


\textbf{Failures and Risk Analysis for AI Scientist Systems.}
There are very few studies that thoroughly analyze or report the risks and failure cases of AI Scientist systems.
Although \cite{tang2024prioritizing} summarizes the risks that AI Scientists may pose, it focuses on hypothesis-based potential risks, rather than empirically observed risks or failures.
While \cite{beel2025evaluating} provides an in-depth analysis of failure cases in AI Scientist-v1~\citep{lu2024ai}, these findings are based on the early AI Scientist, and the analysis was not conducted from the developer's perspective.
While \cite{luo2025more} examines four failure modes (benchmark selection, data leakage, metric misuse, and post-hoc bias), their analysis is limited to experimental diagnostics and early AI Scientists without modern coding. Thus, their analysis remains somewhat limited in scope.

Therefore, we will provide a more comprehensive report on the various risks identified during the development of our state-of-the-art AI Scientist, in order to deepen the community's understanding of AI Scientists.
